\documentclass[11pt,a4paper]{article}

\usepackage[margin=1in]{geometry}
\usepackage{amsmath,amssymb,mathtools}
\usepackage[most]{tcolorbox}
\usepackage{enumitem}
\usepackage{hyperref}
\usepackage{fancyhdr}

\setlength{\parindent}{0pt}
\setlength{\parskip}{0.6em}

% ---------- Page style ----------
\pagestyle{fancy}
\fancyhf{}
\lhead{Advanced Statistical Mechanics / Statistical Mechanics I}
\rhead{Solution Manual I}
\cfoot{\thepage}

% ---------- Hyperref ----------
\hypersetup{
	colorlinks=true,
	linkcolor=blue,
	urlcolor=blue,
	citecolor=blue
}

% ---------- Counters ----------
\newcounter{problem}
\newcounter{saveequation}

% Rounded box for each problem, with automatic numbering
\newtcolorbox{problembox}{
	enhanced,
	breakable,
	colback=white,
	colframe=black,
	boxrule=0.8pt,
	arc=2.5mm,
	left=2mm,
	right=2mm,
	top=1.5mm,
	bottom=1.5mm,
	before upper={\refstepcounter{problem}\textbf{(\theproblem)}\quad\ignorespaces}
}

% Plain solution environment
% Equations inside are numbered as problem.equation, e.g. 3.1, 3.2
% Equation numbering outside remains unchanged
\newenvironment{solution}
{
	\par
	\setcounter{saveequation}{\value{equation}}
	\begingroup
	\setcounter{equation}{0}
	\renewcommand{\theequation}{\arabic{problem}.\arabic{equation}}
	\renewcommand{\theHequation}{sol.\arabic{problem}.\arabic{equation}}
	\noindent\textbf{Solution.}\quad
}
{
	\par
	\endgroup
	\setcounter{equation}{\value{saveequation}}
}

% Background matter box
\definecolor{bgmatterback}{RGB}{245,248,252}
\definecolor{bgmatterframe}{RGB}{120,140,170}
\definecolor{bgmattertitle}{RGB}{70,92,125}

\newtcolorbox{backgroundmatter}[1][]{
	enhanced,
	breakable,
	colback=bgmatterback,
	colframe=bgmatterframe,
	coltitle=white,
	colbacktitle=bgmattertitle,
	boxrule=0.9pt,
	arc=2.8mm,
	left=2.5mm,
	right=2.5mm,
	top=1.5mm,
	bottom=1.5mm,
	title=Background Matter,
	fonttitle=\bfseries,
	attach boxed title to top left={xshift=2mm,yshift=-2mm},
	boxed title style={
		arc=2mm,
		boxrule=0pt,
		left=1.5mm,
		right=1.5mm,
		top=0.8mm,
		bottom=0.8mm
	},
	before upper={
		\setlength{\parindent}{2em}
		\setlength{\parskip}{0pt}
	},
	#1
}

% Final answer box
\definecolor{finalansback}{RGB}{248,251,247}
\definecolor{finalansframe}{RGB}{88,122,96}
\definecolor{finalanstitle}{RGB}{63,96,70}

\newtcolorbox{finalanswer}[1][]{
	enhanced,
	breakable,
	colback=finalansback,
	colframe=finalansframe,
	coltitle=white,
	colbacktitle=finalanstitle,
	boxrule=1pt,
	arc=2.8mm,
	left=2.5mm,
	right=2.5mm,
	top=1.5mm,
	bottom=1.5mm,
	title=Final Answer,
	fonttitle=\bfseries,
	attach boxed title to top left={xshift=2mm,yshift=-2mm},
	boxed title style={
		arc=2mm,
		boxrule=0pt,
		left=1.5mm,
		right=1.5mm,
		top=0.8mm,
		bottom=0.8mm
	},
	#1
}

\title{\textbf{Report I}\\[0.4em]
	\large Advanced Statistical Mechanics / Statistical Mechanics I}
\author{Bufan Zheng\\[0.2em]\normalsize Student ID: 35-256419}
\date{\today}

\begin{document}
	
	\maketitle
	
	The final results can be found in \eqref{fin:1}, \eqref{fin:2}, \eqref{fin:3}, \eqref{fin:4}, \eqref{fin:5}, \eqref{fin:6} and \eqref{fin:7}.
	
	\section*{Landau theory}
	
	Let us consider the Landau theory:
	\begin{equation}
		\label{eq:1}
		f = t\phi^2 + u\phi^4 - h\phi
		\qquad
		(t,h \in \mathbb{R},\ u > 0).
		\tag{1}
	\end{equation}
	
	\begin{problembox}
		\label{prob:1}
		In the absence of the magnetic field (i.e., $h = 0$), determine the critical point $t = t_{\mathrm c}$.
	\end{problembox}
	\begin{solution}
		The equilibrium value of the order parameter $\phi$ is determined by the minimum of the free energy, namely:
		\begin{equation}
			\label{eq:1.1}
			\left.\frac{\partial f}{\partial \phi}\right|_{\phi=\phi_{\mathrm{eq}}}=0
		\end{equation}
		
		Setting the external field to zero, $h=0$, we have
		\begin{equation}
			f(\phi)=t\phi^2+u\phi^4.
		\end{equation}
		
		The extremum condition is
		\begin{equation}
			\frac{\partial f}{\partial \phi}=2t\phi+4u\phi^3=2\phi\left(t+2u\phi^2\right)=0.
		\end{equation}
		
		We then obtain the following solutions
		\begin{equation}
			\label{eq:1.4}
			\phi=0,\qquad \phi^2=-\frac{t}{2u}.
		\end{equation}
		
		Using the second derivative
		\begin{equation}
		\frac{\partial^2 f}{\partial \phi^2}=2t+12u\phi^2
		\end{equation}
		we discuss the above solutions as follows:
		
		\begin{itemize}
			\item When $t>0$, only $\phi=0$ is a stable minimum.
			\item When $t<0$, the following two stable minima appear:
		\end{itemize}
		
		\begin{equation}
			\phi=\pm\sqrt{-\frac{t}{2u}}
		\end{equation}
		\begin{finalanswer}
		Therefore, the critical point lies at the intersection of the two kinds of behavior,
			\begin{equation}
				\label{fin:1}
				t_c=0.
			\end{equation}
			\end{finalanswer}
	    \end{solution}
	
	\begin{problembox}
		In the absence of the magnetic field, the specific heat $c$ exhibits the following critical phenomenon around the critical point $t = t_{\mathrm c}$:
		\begin{equation}
			c =
			\begin{cases}
				a_- |t - t_{\mathrm c}|^{-\alpha_-} & (t < t_{\mathrm c}), \\
				a_+ |t - t_{\mathrm c}|^{-\alpha_+} & (t > t_{\mathrm c}).
			\end{cases}
			\tag{2}
		\end{equation}
		Calculate the critical exponents $\alpha_{\pm}$.
	\end{problembox}
	\begin{solution}
		According to the equilibrium order parameter obtained from \eqref{eq:1.4}, the free energy in equilibrium can be determined as:
		
		\begin{itemize}
			\item When $t>0$, the equilibrium point is $\phi=0$, so
			\begin{equation}
				f_{\min}=0.
			\end{equation}
			
			\item When $t<0$, the equilibrium point satisfies
			\begin{equation}
				\phi^2=-\frac{t}{2u}.
			\end{equation}
			Substituting back into the free energy,
			\begin{equation}
				f_{\min}=t\phi^2+u\phi^4=t\left(-\frac{t}{2u}\right)+u\left(\frac{t^2}{4u^2}\right)=-\frac{t^2}{4u}.
			\end{equation}
		\end{itemize}
		
		Therefore
		\begin{equation}
			f_{\min}(t)=
			\begin{cases}
				0, & t>0,\\[6pt]
				-\dfrac{t^2}{4u}, & t<0.
			\end{cases}
		\end{equation}
		
		According to the thermodynamic relation, the specific heat can be expressed by the second derivative of the free energy with respect to temperature as
		\begin{equation}
			c=-\frac{\partial^2 f_{\min}}{\partial t^2}.
		\end{equation}
		
		Thus
		\begin{equation}
			c=
			\begin{cases}
				0, & t>0,\\[6pt]
				\dfrac{1}{2u}, & t<0.
			\end{cases}
		\end{equation}
		
		\begin{finalanswer}
			It can be seen that the specific heat does not diverge near the critical point, but only has a finite jump. This is because what is described by \eqref{eq:1} is actually a system undergoing a second-order phase transition, and
			\begin{equation}
				 \label{fin:2}
				\alpha_{+}=\alpha_{-}=0.
			\end{equation}
		\end{finalanswer}
	\end{solution}
	
	\begin{problembox}
		The spontaneous magnetization $m$ exhibits the following critical phenomenon around the critical point $t = t_{\mathrm c}$:
		\begin{equation}
			m \propto |t - t_{\mathrm c}|^{\beta}
			\qquad
			(t < t_{\mathrm c}).
			\tag{3}
		\end{equation}
		Calculate the critical exponent $\beta$.
	\end{problembox}
	\begin{solution}
		The spontaneous magnetization $m$ is in fact the value of the order parameter $\phi$ at equilibrium. According to the solutions of the previous two problems, when $h=0$ and $t<t_c=0$,
		\begin{equation}
			m=\phi_{\mathrm{eq}}=\pm\sqrt{-\frac{t}{2u}}.
		\end{equation}
		
		Therefore, its magnitude satisfies
		\begin{equation}
			|m|\propto |t-t_c|^{1/2}.
		\end{equation}
		
		With the definition
		\begin{equation}
			m\propto |t-t_c|^{\beta} \qquad (t<t_c)
		\end{equation}
		
		\begin{finalanswer}
			by comparison, we obtain
			\begin{equation}
				\label{fin:3}
				\beta=\frac{1}{2}.
			\end{equation}
		\end{finalanswer}
	\end{solution}
	
	\begin{problembox}
		The zero-field susceptibility $\chi$ exhibits the following critical phenomenon around the critical point $t = t_{\mathrm c}$:
		\begin{equation}
			\chi =
			\begin{cases}
				c_- |t - t_{\mathrm c}|^{-\gamma_-} & (t < t_{\mathrm c}), \\
				c_+ |t - t_{\mathrm c}|^{-\gamma_+} & (t > t_{\mathrm c}).
			\end{cases}
			\tag{4}
		\end{equation}
		Calculate the critical exponents $\gamma_{\pm}$ and the ratio $c_+/c_-$ of the coefficients.
		
		\medskip
		\textbf{[Note]} While the coefficients $c_{\pm}$ themselves are not universal, their ratio $c_+/c_-$ is universal.
	\end{problembox}
	\begin{solution}
		Now we include the external field and consider the order parameter in equilibrium. Following \eqref{eq:1.1}, we write down the following equation. In order to keep the notation consistent with the later calculation of the magnetic susceptibility, $\phi_{\mathrm{eq}}$ has already been replaced in advance by $m$ here:
		\begin{equation}
			\label{eq:4.1}
			h=2tm+4um^3.
		\end{equation}
		
		According to the definition of the magnetic susceptibility,
		\begin{equation}
			\chi=\left(\frac{\partial m}{\partial h}\right)_{h\to 0}
			=\left(\frac{\partial h}{\partial m}\right)^{-1}_{h\to 0}
			=\left.\left(2t+12um^2\right)^{-1}\right|_{h\to 0}.
		\end{equation}
		
		Now, using the discussion in Problem \ref{prob:1}, we discuss the two cases as follows:
		\begin{itemize}
			\item When $t>0$, the zero-field equilibrium solution is $m=0$, so
			\begin{equation}
				\chi_{+}=\frac{1}{2t}=\frac{1}{2}|t-t_c|^{-1}.
			\end{equation}
			
			Therefore
			\begin{equation}
				\gamma_{+}=1,\qquad c_{+}=\frac{1}{2}.
			\end{equation}
			
			\item When $t<0$, the zero-field equilibrium solution satisfies
			\begin{equation}
				m^2=-\frac{t}{2u}.
			\end{equation}
			
			Substituting this in, we obtain
			\begin{equation}
				\chi_{-}^{-1}=2t+12u\left(-\frac{t}{2u}\right)=2t-6t=-4t=4|t|.
			\end{equation}
			
			Therefore
			\begin{equation}
				\chi_{-}=\frac{1}{4|t|}=\frac{1}{4}|t-t_c|^{-1}.
			\end{equation}
			
			by comparison, we obtain
			\begin{equation}
				\gamma_{-}=1,\qquad c_{-}=\frac{1}{4}.
			\end{equation}
		\end{itemize}
		\begin{finalanswer}
			Thus, the critical exponents $\gamma_{\pm}$ and the ratio ${c_{+}}/{c_{-}}$ of the coefficients are:
			\begin{equation}
				\label{fin:4}
				\gamma_\pm =1,\quad \frac{c_{+}}{c_{-}}=\frac{1/2}{1/4}=2.
			\end{equation}
		\end{finalanswer}
	\end{solution}
	
	\begin{problembox}
		At the critical point $t = t_{\mathrm c}$, the magnetization $m$ exhibits the following critical phenomenon:
		\begin{equation}
			m \propto |h|^{1/\delta}.
			\tag{5}
		\end{equation}
		Calculate the critical exponent $\delta$.
	\end{problembox}
	\begin{solution}
		At the critical point $t=t_c=0$, \eqref{eq:4.1} becomes
		\begin{equation}
			h=4um^3.
		\end{equation}
		
		Therefore
		\begin{equation}
			m=\left(\frac{h}{4u}\right)^{1/3}\propto |h|^{1/3}.
		\end{equation}
		
	\begin{finalanswer}
		Comparing with the definition, we obtain
		\begin{equation}
			\label{fin:5}
			\delta=3.
		\end{equation}
	\end{finalanswer}
	\end{solution}
	
	\section*{Optional assignment}
	\begin{backgroundmatter}
		While we have implicitly assumed that $\phi$ is real, allowing $\phi$ to take complex values also provides insight into phase transitions and critical phenomena [C. N. Yang and T. D. Lee, \href{https://doi.org/10.1103/PhysRev.87.404}{Phys. Rev. 87, 404 (1952)}; T. D. Lee and C. N. Yang, \href{https://doi.org/10.1103/PhysRev.87.410}{Phys. Rev. 87, 410 (1952)}]. Specifically, in contrast with the real-valued regime $\phi \in \mathbb{R}$, the Landau theory $f$ in Eq. (1) generally has three stationary points for the complex-valued regime $\phi \in \mathbb{C}$. For certain parameter choices, two of these stationary points coalesce, signaling a phase transition.
		
		For clarity, let us here assume $t \ge 0$ and $h = i\gamma$ $(\gamma \ge 0)$. Then, the stationary solutions for $\phi$ (i.e., magnetization $m$) are generally purely imaginary. As described above, two stationary points merge at $\gamma = \gamma_{\mathrm c}$, around which unique critical phenomena arise, the Yang-Lee edge singularity [M. E. Fisher, \href{https://doi.org/10.1103/PhysRevLett.40.1610}{Phys. Rev. Lett. 40, 1610 (1978)}].
	\end{backgroundmatter}
	
	\begin{problembox}
		(optional) Determine $\gamma_{\mathrm c}$ and obtain the following critical exponent $\Delta$:
		\begin{equation}
			\gamma_{\mathrm c} \propto t^{\Delta}.
			\tag{6}
		\end{equation}
	\end{problembox}
	\begin{solution}
		Setting $\frac{\partial f}{\partial \phi}=0$, we obtain the following equation:
		\begin{equation}
			\label{eq:6.1}
			4u\phi^3+2t\phi-i\gamma=0
		\end{equation}
		
		Using the variable substitution $\phi\to i\varphi$, this can be rewritten as the following cubic equation with real coefficients:
		\begin{equation}
			4u\varphi^3-2t\varphi+\gamma=0
		\end{equation}
		
		Putting it into the standard form:
		\begin{equation}
			\varphi^3+p\varphi+q=0,\qquad p=-\frac{t}{2u},\ q=\frac{\gamma}{4u}
		\end{equation}
		
		Using the Cardano discriminant:
		\begin{equation}
			\Delta_c=\left(\frac{q}{2}\right)^2+\left(\frac{p}{3}\right)^3
		\end{equation}
		
		\begin{itemize}
			\item $\Delta_c>0$, one real root and two complex conjugate roots
			\item $\Delta_c=0$, multiple roots
			\item $\Delta_c<0$, all three roots are real
		\end{itemize}
		
		A multiple root is a second-order zero point, so in addition to the first derivative being zero, it also satisfies that the second derivative is zero:
		\begin{equation}
			\label{eq:6.5}
			\frac{\partial^2 f}{\partial \phi^2}=2t+12u\phi^2=0
		\end{equation}
		
		This property will be used in the discussion of the next problem.
		
	\begin{finalanswer}
			Therefore, the condition for having multiple roots is:
		\begin{equation}
			\label{fin:6}
			\gamma_c=\sqrt{\frac{8t^3}{27u}}\sim t^{3/2}
		\end{equation}
		
		Therefore:
		\begin{equation}
			\Delta=\frac{3}{2}
		\end{equation}
	\end{finalanswer}
	\end{solution}
	
	\begin{problembox}
		(optional) Around the critical point $\gamma = \gamma_{\mathrm c}$ for given $t > 0$, the (imaginary) magnetization $m$ exhibits the following critical phenomenon:
		\begin{equation}
			|m - m_{\mathrm c}| \propto |\gamma - \gamma_{\mathrm c}|^{\sigma}.
			\tag{7}
		\end{equation}
		Calculate the critical exponent $\sigma$.
	\end{problembox}
	\begin{solution}
		Perturb around the equilibrium solution
		\begin{equation}
			\phi=\phi_c+\delta\phi,\qquad \gamma=\gamma_c+\delta\gamma.
		\end{equation}
		
		Substitute this into the equilibrium equation of state \eqref{eq:6.1}
		\begin{equation}
			2t(\phi_c+\delta\phi)+4u(\phi_c+\delta\phi)^3-i(\gamma_c+\delta\gamma)=0.
		\end{equation}
		
		Subtract the equation satisfied by the critical point itself, \eqref{eq:6.1}
		\begin{equation}
			2t\phi_c+4u\phi_c^3-i\gamma_c=0,
		\end{equation}
		
		we obtain
		\begin{equation}
			(2t+12u\phi_c^2)\delta\phi+12u\phi_c(\delta\phi)^2+4u(\delta\phi)^3-i\delta\gamma=0.
		\end{equation}
		
		The imaginary magnetization we want to discuss corresponds to the double root, and the double root satisfies \eqref{eq:6.5}
		\begin{equation}
			2t+12u\phi_c^2=0,
		\end{equation}
		
		so the linear term disappears, leaving only
		\begin{equation}
			12u\phi_c(\delta\phi)^2+4u(\delta\phi)^3-i\delta\gamma=0.
		\end{equation}
		
		When we are sufficiently close to the critical point, the quadratic term is dominant, therefore
		\begin{equation}
			\delta\gamma\propto(\delta\phi)^2.
		\end{equation}
		
		Thus
		\begin{equation}
			|m-m_c|=|\phi-\phi_c|\propto|\gamma-\gamma_c|^{1/2}.
		\end{equation}
		
		\begin{finalanswer}
			Therefore
		\begin{equation}
			\label{fin:7}
			\sigma=\frac{1}{2}.
		\end{equation}
		\end{finalanswer}
	\end{solution}
\end{document}